%!TEX root = ../hutbthesis_main.tex

\chapter{讨论与局限性分析}

\section{系统整体性能综合讨论}

本课题以Sapienza/INSA大学DeformableSimulation开源项目为基础进行改造，创建出适合于无硬件环境的可变形物体实时触摸建模和仿真系统，以肝脏、肾脏、皮肤等腹腔器官为仿真对象，用键盘控制虚拟手进行触觉交互仿真，经过第五章实验验证，系统整体达到了预期的改造目标和功能需求，在算法层实现了完整的触觉建模和可变形仿真，突出了改造工作实际的价值。

从物理仿真性能上来说，系统使用了MuJoCo 3.0物理引擎以及flexcomp柔性组件，用mass-spring FEM网络建模腹腔器官的可变形性，用Newton迭代求解器来保证接触检测和形变计算的准确性，实验中肝脏、肾脏、皮肤三种软组织的形变响应符合生物力学特性，视觉渲染流畅，仿真帧率稳定，验证了MuJoCo引擎在生物组织接触仿真中的优势，相比于Bullet、PhysX等引擎，它的数值稳定性更好，更适合进行精确的软体FEM仿真，这也是本系统选择MuJoCo作为主要引擎的原因。

在触觉反馈算法上，系统把原开源项目里和硬件相耦合的力反馈逻辑，拆解成独立的haptic\_feedback.py模块，完成力信号滤波（ForceFilter）、自适应力计算（AdaptiveForceCalculator）、纹理映射（TextureMapper）三个主要功能，实验结果显示，该算法模块具有较好的独立性及有效性，可以对不同的腹腔器官进行力反馈和纹理区分，即使没有WEART触觉手套硬件支持，也能完成模拟器内部完整的力信号建模，体现出本改造工作的新颖之处。

交互控制以及系统稳定方面，系统依靠硬件解耦改造，达成键盘交互的降级模式，依靠pynput跨平台键盘监听功能，达成虚拟手基本控制，借助改进阻尼最小二乘法（DLS-IK），改良虚拟手运动控制，明显削减了穿模可能性。实验过程中系统长时间运行的CPU、内存占用率都在合理范围内，异常发生率低，证明硬件解耦、算法优化后的系统可靠；但是也发现了键盘离散输入造成虚拟手瞬移式驱动，容易和软体产生非物理穿透，导致FEM求解器发散，出现虚拟手飞出视野、器官爆裂等异常，这与原开源项目VR设计初衷有关，是降级交互的固有局限。

综上所述，系统在没有硬件支持的情况下，实现了原开源项目的全部核心功能，完成了算法层触觉建模和可变形仿真全部要求，改造后的系统具有较好的实用性和可扩展性，可以为腹腔手术触觉训练相关研究提供标准化的实验平台，也为之后接入VR/WEART硬件、实现完整的人机闭环交互打下了坚实的基础。

\section{系统优势与创新点}

\begin{enumerate}
\def\labelenumi{\arabic{enumi}.}
\item
  系统改造适配优势
\end{enumerate}

对原开源项目过分依赖VR（Oculus Rift S）和WEART触觉手套硬件的情况进行了硬件解耦改造，使系统可以在没有硬件的情况下正常工作。因此使用pynput键盘监听模块代替OpenXR VR手部追踪程序，设计TUI键盘交互界面（guis.py）代替原项目中的VR交互，实现虚拟手位移、按压、滑动等基本交互操作，并且修复了原项目XML配置文件的兼容性问题，添加了中文化GUI说明，提高系统可操作性及易用性。这样改造不但没有削弱原项目的使用性能，反而丰富了系统的使用场景，在一般的台式机上可以完成实验或者应用的需要。

\begin{enumerate}
\def\labelenumi{\arabic{enumi}.}
\setcounter{enumi}{1}
\item
  触觉反馈算法创新
\end{enumerate}

本系统最核心的创新之处在于把原来开源项目里同硬件关联的力反馈逻辑，分离成独立的触觉反馈算法层，从而达成算法和硬件的完全解耦，具体创新如下所示

(1)力信号滤波优化，用滑动平均滤波算法（ForceFilter）平滑接触力信号，消除高频噪声和数值抖动，提高触觉反馈的细腻度和稳定性，防止原始信号直接输出造成的触感异常。

(2)自适应力计算机制，设计AdaptiveForceCalculator模块，抛弃传统的简单线性力映射，根据手指穿透器官的深度、接触速度来动态调节反馈力，更符合真实的生物组织力学特性，提高了触觉反馈的真实感。

(3)通过TextureMapper模块创建腹腔器官与触觉纹理的独立映射表，将肝脏、肾脏、皮肤分别映射为VenetianGranite、ProfiledRubberSlow、Smooth等纹理来实现对不同组织的精准识别，并且该设计使之后接入任意触觉硬件无需修改核心算法，只需改变硬件接口即可，极大地提高了系统的可扩展性。

\begin{enumerate}
\def\labelenumi{\arabic{enumi}.}
\setcounter{enumi}{2}
\item
  架构与性能优势
\end{enumerate}

系统沿用原开源项目的分层架构，优化后形成``交互层-GUI→主循环→引擎层→触觉反馈算法层''的完整架构，各层之间通过标准化的接口进行通信，模块耦合度低，便于维护和扩展。引擎层使用了Engine接口（interfaces.py）以及MuJoCoConnector（mujoco\_connector.py）来完成手部mocap驱动、软体接触检测、手指IK求解等任务，保证物理仿真与交互控制的高效协同，四线程并行架构可以提高系统的实时性，在实验中系统的端到端延迟控制在合理范围之内，仿真帧率稳定，可以达到触觉交互的基本需求。此外，系统采用MuJoCo flexcomp元素实现器官建模，通过edge equality constraint（等式约束）+阻尼连接构建mass-spring FEM网络，结合solimp/solref参数定义接触特性，既保证了仿真精度，又兼顾了计算效率，相较于刚体建模（rigid body，仅6个自由度），能够更真实地还原软组织的可变形特性。

\section{系统局限性分析}

\begin{enumerate}
\def\labelenumi{\arabic{enumi}.}
\item
  硬件受限导致的交互体验局限
\end{enumerate}

系统目前处在没有VR/WEART硬件的降级交互模式中，主要缺陷就是缺少了完整的``人在回路（human-in-the-loop）''验证。一方面，没有WEART TouchDIVER触觉手套，无法实现真实的力反馈输出，只能在算法上建模力信号，用户不能通过指尖感受到真实的压力和纹理振动，造成触觉反馈的真实性差、沉浸感低；另一方面，键盘离散输入不能完全取代VR头显的连续手部追踪（mocap），虚拟手的运动呈现为瞬移式，精度降低，容易产生非物理穿透，造成FEM求解器发散，出现虚拟手飞出视野、器官爆裂等异常，这也是实验中出现的主要异常现象。

\begin{enumerate}
\def\labelenumi{\arabic{enumi}.}
\setcounter{enumi}{1}
\item
  物理建模层面的局限
\end{enumerate}

系统用MuJoCo flexcomp元素建立腹腔器官模型，本质上就是基于mass-spring FEM网络的离散化建模，虽然可以满足基本的可变形仿真要求，但是相比高精度有限元模型来说，在复杂的形变以及生物组织力学特性上还有不足。以人体肝脏、肾脏的非线性力学性质，软组织的蠕变、应力松弛等状况来说，目前模型的还原程度较低，而且没有考虑拓扑结构的变化，不能完成腹腔手术时器官切割、撕裂、缝合这些高级动作，从而使得系统在虚拟手术训练方面的深度应用受到限制。另外，MuJoCo的Newton迭代求解器虽然比较稳定，但是对于极端大变形、高刚度接触的情况来说，仍然会存在求解效率降低的现象。

\begin{enumerate}
\def\labelenumi{\arabic{enumi}.}
\setcounter{enumi}{2}
\item
  触觉反馈与算法层面的局限
\end{enumerate}

虽然系统已经完成了算法层的触觉反馈建模，但是由于硬件和算法的设计原因，仍然存在着两个不足，即触觉反馈模态单一，只能实现压力、纹理振动的算法建模，不能模拟温度、剪切力、摩擦系数等更多的感知维度，与真实的生物体触觉体验有差距；纹理渲染的细腻度不够高，目前采用的预设纹理映射模式虽然可以区分出不同的器官，但是与真实的生物组织表面触感还有一定的差别，非常细腻的微纹理辨识度不高。多物体接触检测机制可以区分出手和不同的器官之间的接触，但是手指多部位同时接触同一个器官或者多个器官时，力反馈的分配精度还需要进一步提高。

\begin{enumerate}
\def\labelenumi{\arabic{enumi}.}
\setcounter{enumi}{3}
\item
  系统扩展性与适配性局限
\end{enumerate}

目前系统交互方式只用键盘控制，没有实现双手协同交互、多指精细控制等复杂的操作，不能完全还原腹腔手术时的手部动作，场景规模小，只有肝脏、肾脏、皮肤三种腹腔器官模型，可以交互的物体数量少，在多器官紧密聚集的复杂场景中接触检测的效率和稳定性会受到影响。系统对于硬件的配置有一定的要求，核心算法没有做轻量化的处理，不能很好地适应普通低配置的PC以及中端设备，从而影响了系统的普及和应用。

\section{改进方向与优化策略}

\begin{enumerate}
\def\labelenumi{\arabic{enumi}.}
\item
  接入硬件设备，实现完整闭环交互
\end{enumerate}

这是后面优化的主要方向，主要是为了克服目前硬件所造成的限制。将接入Oculus Rift S VR头显（使用pyopenxr接口）来代替键盘交互，实现手部连续追踪（mocap），解决虚拟手瞬移式驱动的问题，降低FEM求解器发散的概率，提高手部控制精度和沉浸感；同时接入WEART TouchDIVER触觉手套（使用WEART SDK），把算法层的力信号建模结果输出到硬件上，得到真实的压力和纹理振动反馈，完成从视觉、物理仿真到触觉反馈的闭环，满足答辩时对于``真实力反馈''的潜在提问需求。

\begin{enumerate}
\def\labelenumi{\arabic{enumi}.}
\setcounter{enumi}{1}
\item
  优化物理建模，提升仿真精度
\end{enumerate}

为了解决物理建模精度低的问题，后面会从两个方面进行改进，一是采用高精度有限元模型或者数据驱动的非线性材料模型来取代目前的mass-spring FEM网络，提高复杂形变、生物组织力学特性（蠕变、应力松弛）的还原程度，更加符合人体腹腔器官真实的触感；二是增加模型的功能，加入器官切割、撕裂、缝合等拓扑结构变化模块，完善虚拟手术训练场景的主要需求；同时改善MuJoCo求解器的参数，调节solimp/solref接触刚度和恢复系数，提高极端情况下的求解速度和稳定性。

\begin{enumerate}
\def\labelenumi{\arabic{enumi}.}
\setcounter{enumi}{2}
\item
  升级触觉反馈算法，丰富感知体验
\end{enumerate}

在现有的算法基础上，提高触觉反馈的真实性、细腻度，即增加温度反馈、剪切力反馈、滑动摩擦反馈等多种触觉通道，形成更加全面的人体触觉感知体验；二是改进纹理生成算法，抛弃预设的纹理映射模式，依照真实物理摩擦特性来创建动态纹理，加强纹理的细腻程度和真实性；三是改良多物体接触检测以及力反馈分配办法，改善手指多部位接触时的力反馈精确度，防止信号串扰。

\begin{enumerate}
\def\labelenumi{\arabic{enumi}.}
\setcounter{enumi}{3}
\item
  扩展系统功能，提升适配性与实用性
\end{enumerate}

第一，增加交互方式，支持双手协同操作、多指精细控制、手势识别等腹腔手术复杂动作的模拟，还原出腹腔手术手部动作；第二，增加更多的腹腔器官模型（脾脏、胃等），创建复杂的腹腔手术场景，改进多器官聚集场景下接触检测的算法；第三，轻量化核心算法，减少硬件要求，让系统可以在普通PC或者中端设备上运行，提升系统的通用性以及可移植性；第四，丰富数据记录与分析模块，加入更多的定量评价指标，给虚拟手术训练效果评定赋予更为全面的数据支撑。

因此，本系统已经完成了核心改造的目标，实现了在没有硬件环境下进行算法层的触觉建模和可变形仿真的工作，具有很强的创新性和实际意义；根据上面的优化策略可以很好地克服目前存在的不足，提高系统的性能，使该系统更加符合腹腔手术触觉训练的要求，为腹腔手术触觉训练研究与应用提供更好的实验平台。

\textbf{\hfill\break
}
